\section*{Combined (unified) RoBERTa -- results}

Fine-tuned \textbf{RoBERTa} on the \textbf{union of all four datasets} (CYNER + APTNER + ATTACKER + DNRTI) in the \textbf{unified} 9-tag BIO label space -- the transformer counterpart of the paper's BiLSTM ``unified-datasets model'' (Section ``Unified-datasets model'', Table \texttt{tab:bigmodel\_vs\_original}). One model is trained and then evaluated on each dataset's unified-label development set plus the pooled \textbf{combined} development set. Recipe: batch 2, learning rate $2\times10^{-5}$, 10 epochs, linear schedule, warmup 0, seed 42, maximum length 512, bf16. Union-leakage removal was applied (training sentences whose tokens appear in any of the four validation sets were dropped): 18,222 to \textbf{17,086} training sentences (1,136 removed). Training wall time \textbf{4,146 s} (approximately 69 min) on one LUMI MI250X GCD. Source: MLflow experiments \texttt{train-combined-roberta-eval-*} and \texttt{models/combined-roberta/train\_metrics.json}.

\subsection*{Combined-model span-F1 versus the BiLSTM combined model (paper Table \texttt{tab:bigmodel\_vs\_original})}

The paper's table compares a BiLSTM trained on the combined data (``Combined'') against BiLSTMs trained on each single dataset (``Original''). Both BiLSTM columns are reproduced from the paper; \textbf{RoBERTa Combined} is this run, and \textbf{RoBERTa Original} is the single-dataset RoBERTa diagonal reproduced from \texttt{cross\_dataset\_roberta\_results.md} (train and evaluate on the same dataset).

\begin{table}[h!]
\centering
\resizebox{0.9\linewidth}{!}{%
\begin{tabular}{lcccc}
\toprule
\textbf{Val \textbackslash Train} & \textbf{BiLSTM Combined} & \textbf{BiLSTM Original} & \textbf{RoBERTa Combined} & \textbf{RoBERTa Original} \\
\midrule
Combined & 0.38 & - & \textbf{0.53} & - \\
DNRTI    & 0.49 & 0.41 & \textbf{0.52} & 0.63 \\
ATTACKER & 0.33 & 0.23 & \textbf{0.51} & 0.61 \\
APTNER   & 0.30 & 0.41 & \textbf{0.49} & 0.59 \\
CYNER    & 0.37 & 0.39 & \textbf{0.67} & 0.73 \\
\bottomrule
\end{tabular}
}%
\caption{Span-F1 on development datasets given models trained on combined or original datasets.}
\label{tab:bigmodel_vs_original}
\end{table}

The combined RoBERTa model beats the BiLSTM combined model on every evaluation set, most strongly on CYNER (0.37 to 0.67) and ATTACKER (0.33 to 0.51). However, unlike the BiLSTM case (where combined training helped on DNRTI and ATTACKER), the \textbf{single-dataset RoBERTa outperforms the combined RoBERTa on every dataset} (for example, CYNER 0.67 to 0.73 and APTNER 0.49 to 0.59). For a pretrained encoder, pooling the four datasets under unified labels is a net loss relative to per-dataset training, consistent with the paper's conclusion that unification noise dominates the benefit of more data.

\begin{quote}
\textbf{Comparability caveat.} The evaluation sets match exactly: both experiments use the same union-leakage removal (which drops only training sentences, leaving development sets full), so each RoBERTa Combined / RoBERTa Original pair is scored on the identical development tokens, unified gold spans, and \texttt{span\_f1}. The two RoBERTa columns do differ in \textit{training-time label handling}, though: RoBERTa Combined is trained natively on the unified label space, whereas RoBERTa Original is trained on each dataset's original labels and collapsed to the four coarse classes by probability-sum late merge after inference. This is an extra confound not present in the paper's BiLSTM columns (both unified-trained), so treat RoBERTa Original as a strong single-dataset reference rather than a perfectly clean ablation.
\end{quote}

\subsection*{Full metric breakdown -- combined RoBERTa}

Strict span-F1 with its precision and recall, plus the unlabelled and loose (partial-overlap, same-label) span-F1 variants, analogous to the paper's cross-dataset metric tables (\texttt{tab:cross\_eval\_matrix\_precision}, \texttt{\_recall}, \texttt{tab:unlabelled-span-f1-cross-dataset}, \texttt{tab:loose-span-f1-cross-dataset}).


\begin{table}[h!]
\centering
\resizebox{0.9\linewidth}{!}{%
\begin{tabular}{lccccc}
\toprule
\textbf{Eval set} & \textbf{Precision} & \textbf{Recall} & \textbf{Span-F1} & \textbf{Unlabelled F1} & \textbf{Loose F1} \\
\midrule
Combined & 0.53 & 0.53 & 0.53 & 0.64 & 0.60 \\
DNRTI    & 0.60 & 0.46 & 0.52 & 0.63 & 0.58 \\
ATTACKER & 0.56 & 0.48 & 0.51 & 0.59 & 0.59 \\
APTNER   & 0.41 & 0.60 & 0.49 & 0.59 & 0.57 \\
CYNER    & 0.68 & 0.66 & 0.67 & 0.80 & 0.71 \\
\bottomrule
\end{tabular}
}%
\caption{Full metric breakdown for the combined RoBERTa model.}
\label{tab:combined_roberta_metrics}
\end{table}

CYNER is strongest across every metric. DNRTI and ATTACKER are precision-leaning (precision greater than recall), whereas APTNER is the one recall-leaning set (0.60 versus 0.41): the model over-predicts entities on APTNER, consistent with the definition-discrepancy trends discussed in the paper's analysis. The gap between strict and unlabelled F1 (for example, CYNER 0.67 to 0.80) reflects residual label-assignment errors on otherwise correctly delimited spans.
